\chapter{重构计划}

% === 源文件: refactor/clawnet.md ===
\section{Clawnet 重构（协议 + 认证统一）}


\subsection{嗨}


嗨 Peter — 方向很好；这将解锁更简单的用户体验 + 更强的安全性。

\subsection{目的}


单一、严谨的文档用于：

\begin{itemize}
  \item 当前状态：协议、流程、信任边界。
  \item 痛点：审批、多跳路由、UI 重复。
  \item 提议的新状态：一个协议、作用域角色、统一的认证/配对、TLS 固定。
  \item 身份模型：稳定 ID + 可爱的别名。
  \item 迁移计划、风险、开放问题。
\end{itemize}


\subsection{目标（来自讨论）}


\begin{itemize}
  \item 所有客户端使用一个协议（mac 应用、CLI、iOS、Android、无头节点）。
  \item 每个网络参与者都经过认证 + 配对。
  \item 角色清晰：节点 vs 操作者。
  \item 中央审批路由到用户所在位置。
  \item 所有远程流量使用 TLS 加密 + 可选固定。
  \item 最小化代码重复。
  \item 单台机器应该只显示一次（无 UI/节点重复条目）。
\end{itemize}


\subsection{非目标（明确）}


\begin{itemize}
  \item 移除能力分离（仍需要最小权限）。
  \item 不经作用域检查就暴露完整的 Gateway 网关控制平面。
  \item 使认证依赖于人类标签（别名仍然是非安全性的）。
\end{itemize}


\hrulefill


\section{当前状态（现状）}


\subsection{两个协议}


\subsubsection{1) Gateway 网关 WebSocket（控制平面）}


\begin{itemize}
  \item 完整 API 表面：配置、渠道、模型、会话、智能体运行、日志、节点等。
  \item 默认绑定：loopback。通过 SSH/Tailscale 远程访问。
  \item 认证：通过 \texttt{connect} 的令牌/密码。
  \item 无 TLS 固定（依赖 loopback/隧道）。
  \item 代码：
  \item \texttt{src/gateway/server/ws-connection/message-handler.ts}
  \item \texttt{src/gateway/client.ts}
  \item \texttt{docs/gateway/protocol.md}
\end{itemize}


\subsubsection{2) Bridge（节点传输）}


\begin{itemize}
  \item 窄允许列表表面，节点身份 + 配对。
  \item TCP 上的 JSONL；可选 TLS + 证书指纹固定。
  \item TLS 在设备发现 TXT 中公布指纹。
  \item 代码：
  \item \texttt{src/infra/bridge/server/connection.ts}
  \item \texttt{src/gateway/server-bridge.ts}
  \item \texttt{src/node-host/bridge-client.ts}
  \item \texttt{docs/gateway/bridge-protocol.md}
\end{itemize}


\subsection{当前的控制平面客户端}


\begin{itemize}
  \item CLI → 通过 \texttt{callGateway}（\texttt{src/gateway/call.ts}）连接 Gateway 网关 WS。
  \item macOS 应用 UI → Gateway 网关 WS（\texttt{GatewayConnection}）。
  \item Web 控制 UI → Gateway 网关 WS。
  \item ACP → Gateway 网关 WS。
  \item 浏览器控制使用自己的 HTTP 控制服务器。
\end{itemize}


\subsection{当前的节点}


\begin{itemize}
  \item macOS 应用在节点模式下连接到 Gateway 网关 bridge（\texttt{MacNodeBridgeSession}）。
  \item iOS/Android 应用连接到 Gateway 网关 bridge。
  \item 配对 + 每节点令牌存储在 Gateway 网关上。
\end{itemize}


\subsection{当前审批流程（exec）}


\begin{itemize}
  \item 智能体通过 Gateway 网关使用 \texttt{system.run}。
  \item Gateway 网关通过 bridge 调用节点。
  \item 节点运行时决定审批。
  \item UI 提示由 mac 应用显示（当节点 == mac 应用时）。
  \item 节点向 Gateway 网关返回 \texttt{invoke-res}。
  \item 多跳，UI 绑定到节点主机。
\end{itemize}


\subsection{当前的在线状态 + 身份}


\begin{itemize}
  \item 来自 WS 客户端的 Gateway 网关在线状态条目。
  \item 来自 bridge 的节点在线状态条目。
  \item mac 应用可能为同一台机器显示两个条目（UI + 节点）。
  \item 节点身份存储在配对存储中；UI 身份是分开的。
\end{itemize}


\hrulefill


\section{问题/痛点}


\begin{itemize}
  \item 需要维护两个协议栈（WS + Bridge）。
  \item 远程节点上的审批：提示出现在节点主机上，而不是用户所在位置。
  \item TLS 固定仅存在于 bridge；WS 依赖 SSH/Tailscale。
  \item 身份重复：同一台机器显示为多个实例。
  \item 角色模糊：UI + 节点 + CLI 能力没有明确分离。
\end{itemize}


\hrulefill


\section{提议的新状态（Clawnet）}


\subsection{一个协议，两个角色}


带有角色 + 作用域的单一 WS 协议。

\begin{itemize}
  \item \textbf{角色：node}（能力宿主）
  \item \textbf{角色：operator}（控制平面）
  \item 操作者的可选\textbf{作用域}：
  \item \texttt{operator.read}（状态 + 查看）
  \item \texttt{operator.write}（智能体运行、发送）
  \item \texttt{operator.admin}（配置、渠道、模型）
\end{itemize}


\subsubsection{角色行为}


\textbf{Node}

\begin{itemize}
  \item 可以注册能力（\texttt{caps}、\texttt{commands}、permissions）。
  \item 可以接收 \texttt{invoke} 命令（\texttt{system.run}、\texttt{camera.\textit{}、\texttt{canvas.}}、\texttt{screen.record} 等）。
  \item 可以发送事件：\texttt{voice.transcript}、\texttt{agent.request}、\texttt{chat.subscribe}。
  \item 不能调用配置/模型/渠道/会话/智能体控制平面 API。
\end{itemize}


\textbf{Operator}

\begin{itemize}
  \item 完整控制平面 API，受作用域限制。
  \item 接收所有审批。
  \item 不直接执行 OS 操作；路由到节点。
\end{itemize}


\subsubsection{关键规则}


角色是按连接的，不是按设备。一个设备可以分别打开两个角色。

\hrulefill


\section{统一认证 + 配对}


\subsection{客户端身份}


每个客户端提供：

\begin{itemize}
  \item \texttt{deviceId}（稳定的，从设备密钥派生）。
  \item \texttt{displayName}（人类名称）。
  \item \texttt{role} + \texttt{scope} + \texttt{caps} + \texttt{commands}。
\end{itemize}


\subsection{配对流程（统一）}


\begin{itemize}
  \item 客户端未认证连接。
  \item Gateway 网关为该 \texttt{deviceId} 创建\textbf{配对请求}。
  \item 操作者收到提示；批准/拒绝。
  \item Gateway 网关颁发绑定到以下内容的凭证：
  \item 设备公钥
  \item 角色
  \item 作用域
  \item 能力/命令
  \item 客户端持久化令牌，重新认证连接。
\end{itemize}


\subsection{设备绑定认证（避免 bearer 令牌重放）}


首选：设备密钥对。

\begin{itemize}
  \item 设备一次性生成密钥对。
  \item \texttt{deviceId = fingerprint(publicKey)}。
  \item Gateway 网关发送 nonce；设备签名；Gateway 网关验证。
  \item 令牌颁发给公钥（所有权证明），而不是字符串。
\end{itemize}


替代方案：

\begin{itemize}
  \item mTLS（客户端证书）：最强，运维复杂度更高。
  \item 短期 bearer 令牌仅作为临时阶段（早期轮换 + 撤销）。
\end{itemize}


\subsection{静默批准（SSH 启发式）}


精确定义以避免薄弱环节。优选其一：

\begin{itemize}
  \item \textbf{仅限本地}：当客户端通过 loopback/Unix socket 连接时自动配对。
  \item \textbf{通过 SSH 质询}：Gateway 网关颁发 nonce；客户端通过获取它来证明 SSH。
  \item \textbf{物理存在窗口}：在 Gateway 网关主机 UI 上本地批准后，允许在短窗口内（例如 10 分钟）自动配对。
\end{itemize}


始终记录 + 记录自动批准。

\hrulefill


\section{TLS 无处不在（开发 + 生产）}


\subsection{复用现有 bridge TLS}


使用当前 TLS 运行时 + 指纹固定：

\begin{itemize}
  \item \texttt{src/infra/bridge/server/tls.ts}
  \item \texttt{src/node-host/bridge-client.ts} 中的指纹验证逻辑
\end{itemize}


\subsection{应用于 WS}


\begin{itemize}
  \item WS 服务器使用相同的证书/密钥 + 指纹支持 TLS。
  \item WS 客户端可以固定指纹（可选）。
  \item 设备发现为所有端点公布 TLS + 指纹。
  \item 设备发现仅是定位器提示；永远不是信任锚。
\end{itemize}


\subsection{为什么}


\begin{itemize}
  \item 减少对 SSH/Tailscale 的机密性依赖。
  \item 默认情况下使远程移动连接安全。
\end{itemize}


\hrulefill


\section{审批重新设计（集中化）}


\subsection{当前}


审批发生在节点主机上（mac 应用节点运行时）。提示出现在节点运行的地方。

\subsection{提议}


审批是 \textbf{Gateway 网关托管的}，UI 传递给操作者客户端。

\subsubsection{新流程}


\begin{enumerate}
  \item Gateway 网关接收 \texttt{system.run} 意图（智能体）。
  \item Gateway 网关创建审批记录：\texttt{approval.requested}。
  \item 操作者 UI 显示提示。
  \item 审批决定发送到 Gateway 网关：\texttt{approval.resolve}。
  \item 如果批准，Gateway 网关调用节点命令。
  \item 节点执行，返回 \texttt{invoke-res}。
\end{enumerate}


\subsubsection{审批语义（加固）}


\begin{itemize}
  \item 广播到所有操作者；只有活跃的 UI 显示模态框（其他显示 toast）。
  \item 先解决者获胜；Gateway 网关拒绝后续解决为已结算。
  \item 默认超时：N 秒后拒绝（例如 60 秒），记录原因。
  \item 解决需要 \texttt{operator.approvals} 作用域。
\end{itemize}


\subsection{好处}


\begin{itemize}
  \item 提示出现在用户所在位置（mac/手机）。
  \item 远程节点的一致审批。
  \item 节点运行时保持无头；无 UI 依赖。
\end{itemize}


\hrulefill


\section{角色清晰示例}


\subsection{iPhone 应用}


\begin{itemize}
  \item \textbf{Node 角色}用于：麦克风、相机、语音聊天、位置、一键通话。
  \item 可选的 \textbf{operator.read} 用于状态和聊天视图。
  \item 可选的 \textbf{operator.write/admin} 仅在明确启用时。
\end{itemize}


\subsection{macOS 应用}


\begin{itemize}
  \item 默认是 Operator 角色（控制 UI）。
  \item 启用"Mac 节点"时是 Node 角色（system.run、屏幕、相机）。
  \item 两个连接使用相同的 deviceId → 合并的 UI 条目。
\end{itemize}


\subsection{CLI}


\begin{itemize}
  \item 始终是 Operator 角色。
  \item 作用域按子命令派生：
  \item \texttt{status}、\texttt{logs} → read
  \item \texttt{agent}、\texttt{message} → write
  \item \texttt{config}、\texttt{channels} → admin
  \item 审批 + 配对 → \texttt{operator.approvals} / \texttt{operator.pairing}
\end{itemize}


\hrulefill


\section{身份 + 别名}


\subsection{稳定 ID}


认证必需；永不改变。
首选：

\begin{itemize}
  \item 密钥对指纹（公钥哈希）。
\end{itemize}


\subsection{可爱别名（龙虾主题）}


仅人类标签。

\begin{itemize}
  \item 示例：\texttt{scarlet-claw}、\texttt{saltwave}、\texttt{mantis-pinch}。
  \item 存储在 Gateway 网关注册表中，可编辑。
  \item 冲突处理：\texttt{-2}、\texttt{-3}。
\end{itemize}


\subsection{UI 分组}


跨角色的相同 \texttt{deviceId} → 单个"实例"行：

\begin{itemize}
  \item 徽章：\texttt{operator}、\texttt{node}。
  \item 显示能力 + 最后在线。
\end{itemize}


\hrulefill


\section{迁移策略}


\subsection{阶段 0：记录 + 对齐}


\begin{itemize}
  \item 发布此文档。
  \item 盘点所有协议调用 + 审批流程。
\end{itemize}


\subsection{阶段 1：向 WS 添加角色/作用域}


\begin{itemize}
  \item 用 \texttt{role}、\texttt{scope}、\texttt{deviceId} 扩展 \texttt{connect} 参数。
  \item 为 node 角色添加允许列表限制。
\end{itemize}


\subsection{阶段 2：Bridge 兼容性}


\begin{itemize}
  \item 保持 bridge 运行。
  \item 并行添加 WS node 支持。
  \item 通过配置标志限制功能。
\end{itemize}


\subsection{阶段 3：中央审批}


\begin{itemize}
  \item 在 WS 中添加审批请求 + 解决事件。
  \item 更新 mac 应用 UI 以提示 + 响应。
  \item 节点运行时停止提示 UI。
\end{itemize}


\subsection{阶段 4：TLS 统一}


\begin{itemize}
  \item 使用 bridge TLS 运行时为 WS 添加 TLS 配置。
  \item 向客户端添加固定。
\end{itemize}


\subsection{阶段 5：弃用 bridge}


\begin{itemize}
  \item 将 iOS/Android/mac 节点迁移到 WS。
  \item 保持 bridge 作为后备；稳定后移除。
\end{itemize}


\subsection{阶段 6：设备绑定认证}


\begin{itemize}
  \item 所有非本地连接都需要基于密钥的身份。
  \item 添加撤销 + 轮换 UI。
\end{itemize}


\hrulefill


\section{安全说明}


\begin{itemize}
  \item 角色/允许列表在 Gateway 网关边界强制执行。
  \item 没有客户端可以在没有 operator 作用域的情况下获得"完整"API。
  \item \textit{所有}连接都需要配对。
  \item TLS + 固定减少移动设备的 MITM 风险。
  \item SSH 静默批准是便利措施；仍然记录 + 可撤销。
  \item 设备发现永远不是信任锚。
  \item 能力声明通过按平台/类型的服务器允许列表验证。
\end{itemize}


\section{流式传输 + 大型负载（节点媒体）}


WS 控制平面对于小消息没问题，但节点还做：

\begin{itemize}
  \item 相机剪辑
  \item 屏幕录制
  \item 音频流
\end{itemize}


选项：

\begin{enumerate}
  \item WS 二进制帧 + 分块 + 背压规则。
  \item 单独的流式端点（仍然是 TLS + 认证）。
  \item 对于媒体密集型命令保持 bridge 更长时间，最后迁移。
\end{enumerate}


在实现前选择一个以避免漂移。

\section{能力 + 命令策略}


\begin{itemize}
  \item 节点报告的 caps/commands 被视为\textbf{声明}。
  \item Gateway 网关强制执行每平台允许列表。
  \item 任何新命令都需要操作者批准或显式允许列表更改。
  \item 用时间戳审计更改。
\end{itemize}


\section{审计 + 速率限制}


\begin{itemize}
  \item 记录：配对请求、批准/拒绝、令牌颁发/轮换/撤销。
  \item 速率限制配对垃圾和审批提示。
\end{itemize}


\section{协议卫生}


\begin{itemize}
  \item 显式协议版本 + 错误代码。
  \item 重连规则 + 心跳策略。
  \item 在线状态 TTL 和最后在线语义。
\end{itemize}


\hrulefill


\section{开放问题}


\begin{enumerate}
  \item 同时运行两个角色的单个设备：令牌模型
\end{enumerate}

\begin{itemize}
  \item 建议每个角色单独的令牌（node vs operator）。
  \item 相同的 deviceId；不同的作用域；更清晰的撤销。
\end{itemize}


\begin{enumerate}
  \item 操作者作用域粒度
\end{enumerate}

\begin{itemize}
  \item read/write/admin + approvals + pairing（最小可行）。
  \item 以后考虑每功能作用域。
\end{itemize}


\begin{enumerate}
  \item 令牌轮换 + 撤销 UX
\end{enumerate}

\begin{itemize}
  \item 角色更改时自动轮换。
  \item 按 deviceId + 角色撤销的 UI。
\end{itemize}


\begin{enumerate}
  \item 设备发现
\end{enumerate}

\begin{itemize}
  \item 扩展当前 Bonjour TXT 以包含 WS TLS 指纹 + 角色提示。
  \item 仅作为定位器提示处理。
\end{itemize}


\begin{enumerate}
  \item 跨网络审批
\end{enumerate}

\begin{itemize}
  \item 广播到所有操作者客户端；活跃的 UI 显示模态框。
  \item 先响应者获胜；Gateway 网关强制原子性。
\end{itemize}


\hrulefill


\section{总结（TL;DR）}


\begin{itemize}
  \item 当前：WS 控制平面 + Bridge 节点传输。
  \item 痛点：审批 + 重复 + 两个栈。
  \item 提议：一个带有显式角色 + 作用域的 WS 协议，统一配对 + TLS 固定，Gateway 网关托管的审批，稳定设备 ID + 可爱别名。
  \item 结果：更简单的 UX，更强的安全性，更少的重复，更好的移动路由。
\end{itemize}



\clearpage

% === 源文件: refactor/exec-host.md ===
\section{Exec 主机重构计划}


\subsection{目标}


\begin{itemize}
  \item 添加 \texttt{exec.host} + \texttt{exec.security} 以在\textbf{沙箱}、\textbf{Gateway 网关}和\textbf{节点}之间路由执行。
  \item 保持默认\textbf{安全}：除非明确启用，否则不进行跨主机执行。
  \item 将执行拆分为\textbf{无头运行器服务}，通过本地 IPC 连接可选的 UI（macOS 应用）。
  \item 提供\textbf{每智能体}策略、允许列表、询问模式和节点绑定。
  \item 支持\textit{与}或\textit{不与}允许列表一起使用的\textbf{询问模式}。
  \item 跨平台：Unix socket + token 认证（macOS/Linux/Windows 一致性）。
\end{itemize}


\subsection{非目标}


\begin{itemize}
  \item 无遗留允许列表迁移或遗留 schema 支持。
  \item 节点 exec 无 PTY/流式传输（仅聚合输出）。
  \item 除现有 Bridge + Gateway 网关外无新网络层。
\end{itemize}


\subsection{决定（已锁定）}


\begin{itemize}
  \item \textbf{配置键：} \texttt{exec.host} + \texttt{exec.security}（允许每智能体覆盖）。
  \item \textbf{提升：} 保留 \texttt{/elevated} 作为 Gateway 网关完全访问的别名。
  \item \textbf{询问默认：} \texttt{on-miss}。
  \item \textbf{批准存储：} \texttt{\textasciitilde{}/.openclaw/exec-approvals.json}（JSON，无遗留迁移）。
  \item \textbf{运行器：} 无头系统服务；UI 应用托管 Unix socket 用于批准。
  \item \textbf{节点身份：} 使用现有 \texttt{nodeId}。
  \item \textbf{Socket 认证：} Unix socket + token（跨平台）；如需要稍后拆分。
  \item \textbf{节点主机状态：} \texttt{\textasciitilde{}/.openclaw/node.json}（节点 id + 配对 token）。
  \item \textbf{macOS exec 主机：} 在 macOS 应用内运行 \texttt{system.run}；节点主机服务通过本地 IPC 转发请求。
  \item \textbf{无 XPC helper：} 坚持使用 Unix socket + token + 对等检查。
\end{itemize}


\subsection{关键概念}


\subsubsection{主机}


\begin{itemize}
  \item \texttt{sandbox}：Docker exec（当前行为）。
  \item \texttt{gateway}：在 Gateway 网关主机上执行。
  \item \texttt{node}：通过 Bridge 在节点运行器上执行（\texttt{system.run}）。
\end{itemize}


\subsubsection{安全模式}


\begin{itemize}
  \item \texttt{deny}：始终阻止。
  \item \texttt{allowlist}：仅允许匹配项。
  \item \texttt{full}：允许一切（等同于提升模式）。
\end{itemize}


\subsubsection{询问模式}


\begin{itemize}
  \item \texttt{off}：从不询问。
  \item \texttt{on-miss}：仅在允许列表不匹配时询问。
  \item \texttt{always}：每次都询问。
\end{itemize}


询问\textbf{独立于}允许列表；允许列表可与 \texttt{always} 或 \texttt{on-miss} 一起使用。

\subsubsection{策略解析（每次执行）}


\begin{enumerate}
  \item 解析 \texttt{exec.host}（工具参数 → 智能体覆盖 → 全局默认）。
  \item 解析 \texttt{exec.security} 和 \texttt{exec.ask}（相同优先级）。
  \item 如果主机是 \texttt{sandbox}，继续本地沙箱执行。
  \item 如果主机是 \texttt{gateway} 或 \texttt{node}，在该主机上应用安全 + 询问策略。
\end{enumerate}


\subsection{默认安全}


\begin{itemize}
  \item 默认 \texttt{exec.host = sandbox}。
  \item \texttt{gateway} 和 \texttt{node} 默认 \texttt{exec.security = deny}。
  \item 默认 \texttt{exec.ask = on-miss}（仅在安全允许时相关）。
  \item 如果未设置节点绑定，\textbf{智能体可以定向任何节点}，但仅在策略允许时。
\end{itemize}


\subsection{配置表面}


\subsubsection{工具参数}


\begin{itemize}
  \item \texttt{exec.host}（可选）：\texttt{sandbox | gateway | node}。
  \item \texttt{exec.security}（可选）：\texttt{deny | allowlist | full}。
  \item \texttt{exec.ask}（可选）：\texttt{off | on-miss | always}。
  \item \texttt{exec.node}（可选）：当 \texttt{host=node} 时使用的节点 id/名称。
\end{itemize}


\subsubsection{配置键（全局）}


\begin{itemize}
  \item \texttt{tools.exec.host}
  \item \texttt{tools.exec.security}
  \item \texttt{tools.exec.ask}
  \item \texttt{tools.exec.node}（默认节点绑定）
\end{itemize}


\subsubsection{配置键（每智能体）}


\begin{itemize}
  \item \texttt{agents.list[].tools.exec.host}
  \item \texttt{agents.list[].tools.exec.security}
  \item \texttt{agents.list[].tools.exec.ask}
  \item \texttt{agents.list[].tools.exec.node}
\end{itemize}


\subsubsection{别名}


\begin{itemize}
  \item \texttt{/elevated on} = 为智能体会话设置 \texttt{tools.exec.host=gateway}、\texttt{tools.exec.security=full}。
  \item \texttt{/elevated off} = 为智能体会话恢复之前的 exec 设置。
\end{itemize}


\subsection{批准存储（JSON）}


路径：\texttt{\textasciitilde{}/.openclaw/exec-approvals.json}

用途：

\begin{itemize}
  \item \textbf{执行主机}（Gateway 网关或节点运行器）的本地策略 + 允许列表。
  \item 无 UI 可用时的询问回退。
  \item UI 客户端的 IPC 凭证。
\end{itemize}


建议的 schema（v1）：

\begin{lstlisting}[language=json]
{
  "version": 1,
  "socket": {
    "path": "~/.openclaw/exec-approvals.sock",
    "token": "base64-opaque-token"
  },
  "defaults": {
    "security": "deny",
    "ask": "on-miss",
    "askFallback": "deny"
  },
  "agents": {
    "agent-id-1": {
      "security": "allowlist",
      "ask": "on-miss",
      "allowlist": [
        {
          "pattern": "~/Projects/**/bin/rg",
          "lastUsedAt": 0,
          "lastUsedCommand": "rg -n TODO",
          "lastResolvedPath": "/Users/user/Projects/.../bin/rg"
        }
      ]
    }
  }
}
\end{lstlisting}


注意事项：

\begin{itemize}
  \item 无遗留允许列表格式。
  \item \texttt{askFallback} 仅在需要 \texttt{ask} 且无法访问 UI 时应用。
  \item 文件权限：\texttt{0600}。
\end{itemize}


\subsection{运行器服务（无头）}


\subsubsection{角色}


\begin{itemize}
  \item 在本地强制执行 \texttt{exec.security} + \texttt{exec.ask}。
  \item 执行系统命令并返回输出。
  \item 为 exec 生命周期发出 Bridge 事件（可选但推荐）。
\end{itemize}


\subsubsection{服务生命周期}


\begin{itemize}
  \item macOS 上的 Launchd/daemon；Linux/Windows 上的系统服务。
  \item 批准 JSON 是执行主机本地的。
  \item UI 托管本地 Unix socket；运行器按需连接。
\end{itemize}


\subsection{UI 集成（macOS 应用）}


\subsubsection{IPC}


\begin{itemize}
  \item Unix socket 位于 \texttt{\textasciitilde{}/.openclaw/exec-approvals.sock}（0600）。
  \item Token 存储在 \texttt{exec-approvals.json}（0600）中。
  \item 对等检查：仅同 UID。
  \item 挑战/响应：nonce + HMAC(token, request-hash) 防止重放。
  \item 短 TTL（例如 10s）+ 最大负载 + 速率限制。
\end{itemize}


\subsubsection{询问流程（macOS 应用 exec 主机）}


\begin{enumerate}
  \item 节点服务从 Gateway 网关接收 \texttt{system.run}。
  \item 节点服务连接到本地 socket 并发送提示/exec 请求。
  \item 应用验证对等 + token + HMAC + TTL，然后在需要时显示对话框。
  \item 应用在 UI 上下文中执行命令并返回输出。
  \item 节点服务将输出返回给 Gateway 网关。
\end{enumerate}


如果 UI 缺失：

\begin{itemize}
  \item 应用 \texttt{askFallback}（\texttt{deny|allowlist|full}）。
\end{itemize}


\subsubsection{图示（SCI）}


\begin{lstlisting}
Agent -> Gateway -> Bridge -> Node Service (TS)
                         |  IPC (UDS + token + HMAC + TTL)
                         v
                     Mac App (UI + TCC + system.run)
\end{lstlisting}


\subsection{节点身份 + 绑定}


\begin{itemize}
  \item 使用 Bridge 配对中的现有 \texttt{nodeId}。
  \item 绑定模型：
  \item \texttt{tools.exec.node} 将智能体限制为特定节点。
  \item 如果未设置，智能体可以选择任何节点（策略仍强制执行默认值）。
  \item 节点选择解析：
  \item \texttt{nodeId} 精确匹配
  \item \texttt{displayName}（规范化）
  \item \texttt{remoteIp}
  \item \texttt{nodeId} 前缀（>= 6 字符）
\end{itemize}


\subsection{事件}


\subsubsection{谁看到事件}


\begin{itemize}
  \item 系统事件是\textbf{每会话}的，在下一个提示时显示给智能体。
  \item 存储在 Gateway 网关内存队列中（\texttt{enqueueSystemEvent}）。
\end{itemize}


\subsubsection{事件文本}


\begin{itemize}
  \item \texttt{Exec started (node=, id=)}
  \item \texttt{Exec finished (node=, id=, code=)} + 可选输出尾部
  \item \texttt{Exec denied (node=, id=, )}
\end{itemize}


\subsubsection{传输}


选项 A（推荐）：

\begin{itemize}
  \item 运行器发送 Bridge \texttt{event} 帧 \texttt{exec.started} / \texttt{exec.finished}。
  \item Gateway 网关 \texttt{handleBridgeEvent} 将这些映射到 \texttt{enqueueSystemEvent}。
\end{itemize}


选项 B：

\begin{itemize}
  \item Gateway 网关 \texttt{exec} 工具直接处理生命周期（仅同步）。
\end{itemize}


\subsection{Exec 流程}


\subsubsection{沙箱主机}


\begin{itemize}
  \item 现有 \texttt{exec} 行为（Docker 或无沙箱时的主机）。
  \item 仅在非沙箱模式下支持 PTY。
\end{itemize}


\subsubsection{Gateway 网关主机}


\begin{itemize}
  \item Gateway 网关进程在其自己的机器上执行。
  \item 强制执行本地 \texttt{exec-approvals.json}（安全/询问/允许列表）。
\end{itemize}


\subsubsection{节点主机}


\begin{itemize}
  \item Gateway 网关调用 \texttt{node.invoke} 配合 \texttt{system.run}。
  \item 运行器强制执行本地批准。
  \item 运行器返回聚合的 stdout/stderr。
  \item 可选的 Bridge 事件用于开始/完成/拒绝。
\end{itemize}


\subsection{输出上限}


\begin{itemize}
  \item 组合 stdout+stderr 上限为 \textbf{200k}；为事件保留\textbf{尾部 20k}。
  \item 使用清晰的后缀截断（例如 \texttt{"… (truncated)"}）。
\end{itemize}


\subsection{斜杠命令}


\begin{itemize}
  \item \texttt{/exec host= security= ask= node=}
  \item 每智能体、每会话覆盖；除非通过配置保存，否则非持久。
  \item \texttt{/elevated on|off|ask|full} 仍然是 \texttt{host=gateway security=full} 的快捷方式（\texttt{full} 跳过批准）。
\end{itemize}


\subsection{跨平台方案}


\begin{itemize}
  \item 运行器服务是可移植的执行目标。
  \item UI 是可选的；如果缺失，应用 \texttt{askFallback}。
  \item Windows/Linux 支持相同的批准 JSON + socket 协议。
\end{itemize}


\subsection{实现阶段}


\subsubsection{阶段 1：配置 + exec 路由}


\begin{itemize}
  \item 为 \texttt{exec.host}、\texttt{exec.security}、\texttt{exec.ask}、\texttt{exec.node} 添加配置 schema。
  \item 更新工具管道以遵守 \texttt{exec.host}。
  \item 添加 \texttt{/exec} 斜杠命令并保留 \texttt{/elevated} 别名。
\end{itemize}


\subsubsection{阶段 2：批准存储 + Gateway 网关强制执行}


\begin{itemize}
  \item 实现 \texttt{exec-approvals.json} 读取器/写入器。
  \item 为 \texttt{gateway} 主机强制执行允许列表 + 询问模式。
  \item 添加输出上限。
\end{itemize}


\subsubsection{阶段 3：节点运行器强制执行}


\begin{itemize}
  \item 更新节点运行器以强制执行允许列表 + 询问。
  \item 添加 Unix socket 提示桥接到 macOS 应用 UI。
  \item 连接 \texttt{askFallback}。
\end{itemize}


\subsubsection{阶段 4：事件}


\begin{itemize}
  \item 为 exec 生命周期添加节点 → Gateway 网关 Bridge 事件。
  \item 映射到 \texttt{enqueueSystemEvent} 用于智能体提示。
\end{itemize}


\subsubsection{阶段 5：UI 完善}


\begin{itemize}
  \item Mac 应用：允许列表编辑器、每智能体切换器、询问策略 UI。
  \item 节点绑定控制（可选）。
\end{itemize}


\subsection{测试计划}


\begin{itemize}
  \item 单元测试：允许列表匹配（glob + 不区分大小写）。
  \item 单元测试：策略解析优先级（工具参数 → 智能体覆盖 → 全局）。
  \item 集成测试：节点运行器拒绝/允许/询问流程。
  \item Bridge 事件测试：节点事件 → 系统事件路由。
\end{itemize}


\subsection{开放风险}


\begin{itemize}
  \item UI 不可用：确保遵守 \texttt{askFallback}。
  \item 长时间运行的命令：依赖超时 + 输出上限。
  \item 多节点歧义：除非有节点绑定或显式节点参数，否则报错。
\end{itemize}


\subsection{相关文档}


\begin{itemize}
  \item \href{/tools/exec}{Exec 工具}
  \item \href{/tools/exec-approvals}{执行批准}
  \item \href{/nodes}{节点}
  \item \href{/tools/elevated}{提升模式}
\end{itemize}



\clearpage

% === 源文件: refactor/outbound-session-mirroring.md ===
\section{出站会话镜像重构（Issue \#1520）}


\subsection{状态}


\begin{itemize}
  \item 进行中。
  \item 核心 + 插件渠道路由已更新以支持出站镜像。
  \item Gateway 网关发送现在在省略 sessionKey 时派生目标会话。
\end{itemize}


\subsection{背景}


出站发送被镜像到\textit{当前}智能体会话（工具会话键）而不是目标渠道会话。入站路由使用渠道/对等方会话键，因此出站响应落在错误的会话中，首次联系的目标通常缺少会话条目。

\subsection{目标}


\begin{itemize}
  \item 将出站消息镜像到目标渠道会话键。
  \item 在缺失时为出站创建会话条目。
  \item 保持线程/话题作用域与入站会话键对齐。
  \item 涵盖核心渠道加内置扩展。
\end{itemize}


\subsection{实现摘要}


\begin{itemize}
  \item 新的出站会话路由辅助器：
  \item \texttt{src/infra/outbound/outbound-session.ts}
  \item \texttt{resolveOutboundSessionRoute} 使用 \texttt{buildAgentSessionKey}（dmScope + identityLinks）构建目标 sessionKey。
  \item \texttt{ensureOutboundSessionEntry} 通过 \texttt{recordSessionMetaFromInbound} 写入最小的 \texttt{MsgContext}。
  \item \texttt{runMessageAction}（发送）派生目标 sessionKey 并将其传递给 \texttt{executeSendAction} 进行镜像。
  \item \texttt{message-tool} 不再直接镜像；它只从当前会话键解析 agentId。
  \item 插件发送路径使用派生的 sessionKey 通过 \texttt{appendAssistantMessageToSessionTranscript} 进行镜像。
  \item Gateway 网关发送在未提供时派生目标会话键（默认智能体），并确保会话条目。
\end{itemize}


\subsection{线程/话题处理}


\begin{itemize}
  \item Slack：replyTo/threadId -> \texttt{resolveThreadSessionKeys}（后缀）。
  \item Discord：threadId/replyTo -> \texttt{resolveThreadSessionKeys}，\texttt{useSuffix=false} 以匹配入站（线程频道 id 已经作用域会话）。
  \item Telegram：话题 ID 通过 \texttt{buildTelegramGroupPeerId} 映射到 \texttt{chatId:topic:}。
\end{itemize}


\subsection{涵盖的扩展}


\begin{itemize}
  \item Matrix、MS Teams、Mattermost、BlueBubbles、Nextcloud Talk、Zalo、Zalo Personal、Nostr、Tlon。
  \item 注意：
  \item Mattermost 目标现在为私信会话键路由去除 \texttt{@}。
  \item Zalo Personal 对 1:1 目标使用私信对等方类型（仅当存在 \texttt{group:} 时才使用群组）。
  \item BlueBubbles 群组目标去除 \texttt{chat\_*} 前缀以匹配入站会话键。
  \item Slack 自动线程镜像不区分大小写地匹配频道 id。
  \item Gateway 网关发送在镜像前将提供的会话键转换为小写。
\end{itemize}


\subsection{决策}


\begin{itemize}
  \item \textbf{Gateway 网关发送会话派生}：如果提供了 \texttt{sessionKey}，则使用它。如果省略，从目标 + 默认智能体派生 sessionKey 并镜像到那里。
  \item \textbf{会话条目创建}：始终使用 \texttt{recordSessionMetaFromInbound}，\texttt{Provider/From/To/ChatType/AccountId/Originating*} 与入站格式对齐。
  \item \textbf{目标规范化}：出站路由在可用时使用解析后的目标（\texttt{resolveChannelTarget} 之后）。
  \item \textbf{会话键大小写}：在写入和迁移期间将会话键规范化为小写。
\end{itemize}


\subsection{添加/更新的测试}


\begin{itemize}
  \item \texttt{src/infra/outbound/outbound-session.test.ts}
  \item Slack 线程会话键。
  \item Telegram 话题会话键。
  \item dmScope identityLinks 与 Discord。
  \item \texttt{src/agents/tools/message-tool.test.ts}
  \item 从会话键派生 agentId（不传递 sessionKey）。
  \item \texttt{src/gateway/server-methods/send.test.ts}
  \item 在省略时派生会话键并创建会话条目。
\end{itemize}


\subsection{待处理项目 / 后续跟进}


\begin{itemize}
  \item 语音通话插件使用自定义的 \texttt{voice:} 会话键。出站映射在这里没有标准化；如果 message-tool 应该支持语音通话发送，请添加显式映射。
  \item 确认是否有任何外部插件使用内置集之外的非标准 \texttt{From/To} 格式。
\end{itemize}


\subsection{涉及的文件}


\begin{itemize}
  \item \texttt{src/infra/outbound/outbound-session.ts}
  \item \texttt{src/infra/outbound/outbound-send-service.ts}
  \item \texttt{src/infra/outbound/message-action-runner.ts}
  \item \texttt{src/agents/tools/message-tool.ts}
  \item \texttt{src/gateway/server-methods/send.ts}
  \item 测试：
  \item \texttt{src/infra/outbound/outbound-session.test.ts}
  \item \texttt{src/agents/tools/message-tool.test.ts}
  \item \texttt{src/gateway/server-methods/send.test.ts}
\end{itemize}



\clearpage

% === 源文件: refactor/plugin-sdk.md ===
\section{插件 SDK + 运行时重构计划}


目标：每个消息连接器都是一个插件（内置或外部），使用统一稳定的 API。
插件不直接从 \texttt{src/**} 导入任何内容。所有依赖项均通过 SDK 或运行时获取。

\subsection{为什么现在做}


\begin{itemize}
  \item 当前连接器混用多种模式：直接导入核心模块、仅 dist 的桥接方式以及自定义辅助函数。
  \item 这使得升级变得脆弱，并阻碍了干净的外部插件接口。
\end{itemize}


\subsection{目标架构（两层）}


\subsubsection{1）插件 SDK（编译时，稳定，可发布）}


范围：类型、辅助函数和配置工具。无运行时状态，无副作用。

内容（示例）：

\begin{itemize}
  \item 类型：\texttt{ChannelPlugin}、适配器、\texttt{ChannelMeta}、\texttt{ChannelCapabilities}、\texttt{ChannelDirectoryEntry}。
  \item 配置辅助函数：\texttt{buildChannelConfigSchema}、\texttt{setAccountEnabledInConfigSection}、\texttt{deleteAccountFromConfigSection}、
\end{itemize}

\texttt{applyAccountNameToChannelSection}。
\begin{itemize}
  \item 配对辅助函数：\texttt{PAIRING\_APPROVED\_MESSAGE}、\texttt{formatPairingApproveHint}。
  \item 新手引导辅助函数：\texttt{promptChannelAccessConfig}、\texttt{addWildcardAllowFrom}、新手引导类型。
  \item 工具参数辅助函数：\texttt{createActionGate}、\texttt{readStringParam}、\texttt{readNumberParam}、\texttt{readReactionParams}、\texttt{jsonResult}。
  \item 文档链接辅助函数：\texttt{formatDocsLink}。
\end{itemize}


交付方式：

\begin{itemize}
  \item 以 \texttt{openclaw/plugin-sdk} 发布（或从核心以 \texttt{openclaw/plugin-sdk} 导出）。
  \item 使用语义化版本控制，提供明确的稳定性保证。
\end{itemize}


\subsubsection{2）插件运行时（执行层，注入式）}


范围：所有涉及核心运行时行为的内容。
通过 \texttt{OpenClawPluginApi.runtime} 访问，确保插件永远不会导入 \texttt{src/**}。

建议的接口（最小但完整）：

\begin{lstlisting}[language=typescript]
export type PluginRuntime = {
  channel: {
    text: {
      chunkMarkdownText(text: string, limit: number): string[];
      resolveTextChunkLimit(cfg: OpenClawConfig, channel: string, accountId?: string): number;
      hasControlCommand(text: string, cfg: OpenClawConfig): boolean;
    };
    reply: {
      dispatchReplyWithBufferedBlockDispatcher(params: {
        ctx: unknown;
        cfg: unknown;
        dispatcherOptions: {
          deliver: (payload: {
            text?: string;
            mediaUrls?: string[];
            mediaUrl?: string;
          }) => void | Promise<void>;
          onError?: (err: unknown, info: { kind: string }) => void;
        };
      }): Promise<void>;
      createReplyDispatcherWithTyping?: unknown; // adapter for Teams-style flows
    };
    routing: {
      resolveAgentRoute(params: {
        cfg: unknown;
        channel: string;
        accountId: string;
        peer: { kind: RoutePeerKind; id: string };
      }): { sessionKey: string; accountId: string };
    };
    pairing: {
      buildPairingReply(params: { channel: string; idLine: string; code: string }): string;
      readAllowFromStore(channel: string): Promise<string[]>;
      upsertPairingRequest(params: {
        channel: string;
        id: string;
        meta?: { name?: string };
      }): Promise<{ code: string; created: boolean }>;
    };
    media: {
      fetchRemoteMedia(params: { url: string }): Promise<{ buffer: Buffer; contentType?: string }>;
      saveMediaBuffer(
        buffer: Uint8Array,
        contentType: string | undefined,
        direction: "inbound" | "outbound",
        maxBytes: number,
      ): Promise<{ path: string; contentType?: string }>;
    };
    mentions: {
      buildMentionRegexes(cfg: OpenClawConfig, agentId?: string): RegExp[];
      matchesMentionPatterns(text: string, regexes: RegExp[]): boolean;
    };
    groups: {
      resolveGroupPolicy(
        cfg: OpenClawConfig,
        channel: string,
        accountId: string,
        groupId: string,
      ): {
        allowlistEnabled: boolean;
        allowed: boolean;
        groupConfig?: unknown;
        defaultConfig?: unknown;
      };
      resolveRequireMention(
        cfg: OpenClawConfig,
        channel: string,
        accountId: string,
        groupId: string,
        override?: boolean,
      ): boolean;
    };
    debounce: {
      createInboundDebouncer<T>(opts: {
        debounceMs: number;
        buildKey: (v: T) => string | null;
        shouldDebounce: (v: T) => boolean;
        onFlush: (entries: T[]) => Promise<void>;
        onError?: (err: unknown) => void;
      }): { push: (v: T) => void; flush: () => Promise<void> };
      resolveInboundDebounceMs(cfg: OpenClawConfig, channel: string): number;
    };
    commands: {
      resolveCommandAuthorizedFromAuthorizers(params: {
        useAccessGroups: boolean;
        authorizers: Array<{ configured: boolean; allowed: boolean }>;
      }): boolean;
    };
  };
  logging: {
    shouldLogVerbose(): boolean;
    getChildLogger(name: string): PluginLogger;
  };
  state: {
    resolveStateDir(cfg: OpenClawConfig): string;
  };
};
\end{lstlisting}


备注：

\begin{itemize}
  \item 运行时是访问核心行为的唯一方式。
  \item SDK 故意保持小巧和稳定。
  \item 每个运行时方法都映射到现有的核心实现（无重复代码）。
\end{itemize}


\subsection{迁移计划（分阶段，安全）}


\subsubsection{阶段 0：基础搭建}


\begin{itemize}
  \item 引入 \texttt{openclaw/plugin-sdk}。
  \item 在 \texttt{OpenClawPluginApi} 中添加带有上述接口的 \texttt{api.runtime}。
  \item 在过渡期内保留现有导入方式（添加弃用警告）。
\end{itemize}


\subsubsection{阶段 1：桥接清理（低风险）}


\begin{itemize}
  \item 用 \texttt{api.runtime} 替换每个扩展中的 \texttt{core-bridge.ts}。
  \item 优先迁移 BlueBubbles、Zalo、Zalo Personal（已经接近完成）。
  \item 移除重复的桥接代码。
\end{itemize}


\subsubsection{阶段 2：轻度直接导入的插件}


\begin{itemize}
  \item 将 Matrix 迁移到 SDK + 运行时。
  \item 验证新手引导、目录、群组提及逻辑。
\end{itemize}


\subsubsection{阶段 3：重度直接导入的插件}


\begin{itemize}
  \item 迁移 Microsoft Teams（使用运行时辅助函数最多的插件）。
  \item 确保回复/正在输入的语义与当前行为一致。
\end{itemize}


\subsubsection{阶段 4：iMessage 插件化}


\begin{itemize}
  \item 将 iMessage 移入 \texttt{extensions/imessage}。
  \item 用 \texttt{api.runtime} 替换直接的核心调用。
  \item 保持配置键、CLI 行为和文档不变。
\end{itemize}


\subsubsection{阶段 5：强制执行}


\begin{itemize}
  \item 添加 lint 规则 / CI 检查：禁止 \texttt{extensions/\textbf{} 从 \texttt{src/}} 导入。
  \item 添加插件 SDK/版本兼容性检查（运行时 + SDK 语义化版本）。
\end{itemize}


\subsection{兼容性与版本控制}


\begin{itemize}
  \item SDK：语义化版本控制，已发布，变更有文档记录。
  \item 运行时：按核心版本进行版本控制。添加 \texttt{api.runtime.version}。
  \item 插件声明所需的运行时版本范围（例如 \texttt{openclawRuntime: ">=2026.2.0"}）。
\end{itemize}


\subsection{测试策略}


\begin{itemize}
  \item 适配器级单元测试（使用真实核心实现验证运行时函数）。
  \item 每个插件的黄金测试：确保行为无偏差（路由、配对、允许列表、提及过滤）。
  \item CI 中使用单个端到端插件示例（安装 + 运行 + 冒烟测试）。
\end{itemize}


\subsection{待解决问题}


\begin{itemize}
  \item SDK 类型托管在哪里：独立包还是核心导出？
  \item 运行时类型分发：在 SDK 中（仅类型）还是在核心中？
  \item 如何为内置插件与外部插件暴露文档链接？
  \item 过渡期间是否允许仓库内插件有限地直接导入核心模块？
\end{itemize}


\subsection{成功标准}


\begin{itemize}
  \item 所有渠道连接器都是使用 SDK + 运行时的插件。
  \item \texttt{extensions/\textbf{} 不再从 \texttt{src/}} 导入。
  \item 新连接器模板仅依赖 SDK + 运行时。
  \item 外部插件可以在无需访问核心源码的情况下进行开发和更新。
\end{itemize}


相关文档：\href{/tools/plugin}{插件}、\href{/channels/index}{渠道}、\href{/gateway/configuration}{配置}。


\clearpage

% === 源文件: refactor/strict-config.md ===
\section{严格配置验证（仅通过 doctor 进行迁移）}


\subsection{目标}


\begin{itemize}
  \item \textbf{在所有地方拒绝未知配置键}（根级 + 嵌套）。
  \item \textbf{拒绝没有 schema 的插件配置}；不加载该插件。
  \item \textbf{移除加载时的旧版自动迁移}；迁移仅通过 doctor 运行。
  \item \textbf{启动时自动运行 doctor（dry-run）}；如果无效，阻止非诊断命令。
\end{itemize}


\subsection{非目标}


\begin{itemize}
  \item 加载时的向后兼容性（旧版键不会自动迁移）。
  \item 静默丢弃无法识别的键。
\end{itemize}


\subsection{严格验证规则}


\begin{itemize}
  \item 配置必须在每个层级精确匹配 schema。
  \item 未知键是验证错误（根级或嵌套都不允许透传）。
  \item \texttt{plugins.entries..config} 必须由插件的 schema 验证。
  \item 如果插件缺少 schema，\textbf{拒绝插件加载}并显示清晰的错误。
  \item 未知的 \texttt{channels.} 键是错误，除非插件清单声明了该渠道 id。
  \item 所有插件都需要插件清单（\texttt{openclaw.plugin.json}）。
\end{itemize}


\subsection{插件 schema 强制执行}


\begin{itemize}
  \item 每个插件为其配置提供严格的 JSON Schema（内联在清单中）。
  \item 插件加载流程：
\end{itemize}

\begin{enumerate}
  \item 解析插件清单 + schema（\texttt{openclaw.plugin.json}）。
  \item 根据 schema 验证配置。
  \item 如果缺少 schema 或配置无效：阻止插件加载，记录错误。
\end{enumerate}

\begin{itemize}
  \item 错误消息包括：
  \item 插件 id
  \item 原因（缺少 schema / 配置无效）
  \item 验证失败的路径
  \item 禁用的插件保留其配置，但 Doctor + 日志会显示警告。
\end{itemize}


\subsection{Doctor 流程}


\begin{itemize}
  \item 每次加载配置时都会运行 Doctor（默认 dry-run）。
  \item 如果配置无效：
  \item 打印摘要 + 可操作的错误。
  \item 指示：\texttt{openclaw doctor --fix}。
  \item \texttt{openclaw doctor --fix}：
  \item 应用迁移。
  \item 移除未知键。
  \item 写入更新后的配置。
\end{itemize}


\subsection{命令门控（当配置无效时）}


允许的命令（仅诊断）：

\begin{itemize}
  \item \texttt{openclaw doctor}
  \item \texttt{openclaw logs}
  \item \texttt{openclaw health}
  \item \texttt{openclaw help}
  \item \texttt{openclaw status}
  \item \texttt{openclaw gateway status}
\end{itemize}


其他所有命令必须硬失败并显示："Config invalid. Run \texttt{openclaw doctor --fix}."

\subsection{错误用户体验格式}


\begin{itemize}
  \item 单个摘要标题。
  \item 分组部分：
  \item 未知键（完整路径）
  \item 旧版键/需要迁移
  \item 插件加载失败（插件 id + 原因 + 路径）
\end{itemize}


\subsection{实现接触点}


\begin{itemize}
  \item \texttt{src/config/zod-schema.ts}：移除根级透传；所有地方使用严格对象。
  \item \texttt{src/config/zod-schema.providers.ts}：确保严格的渠道 schema。
  \item \texttt{src/config/validation.ts}：未知键时失败；不应用旧版迁移。
  \item \texttt{src/config/io.ts}：移除旧版自动迁移；始终运行 doctor dry-run。
  \item \texttt{src/config/legacy*.ts}：将用法移至仅 doctor。
  \item \texttt{src/plugins/*}：添加 schema 注册表 + 门控。
  \item \texttt{src/cli} 中的 CLI 命令门控。
\end{itemize}


\subsection{测试}


\begin{itemize}
  \item 未知键拒绝（根级 + 嵌套）。
  \item 插件缺少 schema → 插件加载被阻止并显示清晰错误。
  \item 无效配置 → Gateway 网关启动被阻止，诊断命令除外。
  \item Doctor dry-run 自动运行；\texttt{doctor --fix} 写入修正后的配置。
\end{itemize}



\clearpage
